\documentclass[10pt,a4paper]{article}
\usepackage{amsmath,amssymb,amsthm,graphicx}
\usepackage{amsmath,amssymb,amsthm,mathtools,hyperref,cleveref}
\newtheorem{theorem}{Theorem}
\newtheorem{lemma}[theorem]{Lemma}
\newtheorem{proposition}[theorem]{Proposition}
\newtheorem{corollary}[theorem]{Corollary}
\newtheorem{definition}[theorem]{Definition}
\newtheorem{conjecture}[theorem]{Conjecture}
\newtheorem{remark}[theorem]{Remark}
\newcommand{\lean}[1]{}
\newcommand{\proves}[1]{}
\newcommand{\uses}[1]{}
\newcommand{\leanok}{}
\newcommand{\chapter}[1]{\section{#1}}
\usepackage{xurl}
% Put any macro and import needed for the project here.
% This will be used by both the web and print versions of the blueprint.
% This file is not meant to be built. Build src/web.tex or src/print.text instead.

% Letters
\newcommand{\C}{\mathbb{C}}
\newcommand{\Z}{\mathbb{Z}}
\newcommand{\bbc}{\mathbb{C}}
\newcommand{\E}{\mathbb{E}}
\newcommand*{\bbe}{\mathbb{E}}
\newcommand{\F}{\mathbb{F}}
\newcommand{\bbf}{\mathbb{F}}
\newcommand{\bbH}{\mathbb{H}}
\newcommand{\bbP}{\mathbb{P}}
\newcommand{\bbI}{\mathbb{I}}
\newcommand{\bbn}{\mathbb{N}}
\newcommand{\bbq}{\mathbb{Q}}
\newcommand{\bbr}{\mathbb{R}}
\newcommand{\bbt}{\mathbb{T}}
\newcommand{\bbz}{\mathbb{Z}}
\newcommand{\N}{\mathbb{N}}

\newcommand{\lo}[1]{\mathcal{L}{#1}}

% Paired delimiters
\newcommand{\abs}[1]{\left\lvert #1\right\rvert}
\newcommand{\Abs}[1]{\lvert #1 \rvert}
\newcommand{\brac}[1]{\left( #1\right)}
\newcommand{\norm}[1]{\lVert #1\rVert}
\newcommand{\inn}[1]{\left\langle #1 \right\rangle}

% Operators
\DeclareMathOperator{\dist}{dist}

\newcommand{\ind}[1]{1_{#1}}
\providecommand{\tup}[1]{{\vec{#1}}}
\title{Approximate homomorphism form of Polynomial Freiman-Ruzsa}
\author{Terence Tao, PFR project contributors}
\date{github:teorth-pfr:4c56de547b2f33b0bc47b92f3bfbcebc6457a776; Apache-2.0}
\begin{document}
\maketitle
\noindent\textit{Editorial note.} Source-organized author material only. All mathematical source spans are retained unchanged. Blueprint lean/proves/uses metadata is hidden as formatting only; numbering is local to this material wrapper.\par\medskip
\noindent\textit{Editorial note.} Ordinary cited prerequisites and concise author arguments are accepted as such. No independent mathematical certification or public ID allocation is claimed.\par\medskip
\noindent\textit{Editorial note.} Formatting repair after actual visual inspection: xurl allows the original long author URL to wrap within the text block.\par\medskip
\noindent\textit{Allowed standard prerequisites (editorial).}\begin{itemize}
\item Standard finite Shannon entropy identities, conditional entropy and finite abelian group algebra are cited in their original Blueprint formulations.
\item Named prior PFR/entropy or BSG results are allowed prerequisites; their original statements are retained where used, without demanding recursively expanded proofs.
\end{itemize}\medskip
\chapter{Proof of PFR}

\begin{lemma}[Ruzsa covering lemma]
\label{ruz-cov}
\lean{Finset.ruzsa_covering_add}
\leanok
If $A,B$ are finite non-empty subsets of a group $G$, then $A$ can be covered by at most $|A+B|/|B|$ translates of $B-B$.
\end{lemma}
\begin{proof}
\leanok
Cover $A$ greedily by disjoint translates of $B$.
\end{proof}



\begin{proposition} \label{pfr-rho}\lean{rho_PFR_conjecture}\leanok For any random variables $Y_1,Y_2$, there exist a subgroup $H$ such that
  $$ 2\rho(U_H) \leq \rho(Y_1) + \rho(Y_2) + 8 d[Y_1;Y_2].$$
\end{proposition}


\begin{corollary}\label{pfr-9-aux}\lean{better_PFR_conjecture_aux0}\leanok
If $|A+A| \leq K|A|$, then there exists a subgroup $H$ and $t\in G$ such that
$|A \cap (H+t)| \geq K^{-4} \sqrt{|A||H|}$, and $|H|/|A|\in[K^{-8},K^8]$.
\end{corollary}



\begin{corollary}\label{pfr-9-aux'}\lean{better_PFR_conjecture_aux}\leanok
If $|A+A| \leq K|A|$, then there exist a subgroup $H$ and a subset $c$ of $G$
with $A \subseteq c + H$, such that $|c| \leq K^{5} |A|^{1/2}/|H|^{1/2}$ and
$|H|/|A|\in[K^{-8},K^8]$.
\end{corollary}


\begin{lemma}[Cauchy--Schwarz bound]
  \label{cs-bound}
  \uses{energy-def}
  \lean{Finset.card_sq_le_card_mul_addEnergy'}
  \leanok

  If $G$ is a finite additive group and $A,B$ are subsets of $G$, then
  $$ \left( \frac{|\{ (a,a') \in A \times A: a+a' \in B \}|}{|G|^2} \right)^2
     \leq \frac{|B|}{|G|}\, E(A). $$
  The left-hand side is the square of the density of a subset of $G \times G$, and the
  right-hand side involves the density of a subset of $G$.
\end{lemma}
\begin{proof}
  \leanok

  If $B$ is empty then the claim is trivial (with the Lean convention $0/0$),
  so without loss of generality $B$ is non-empty.  We can rewrite
  $$ |\{ (a,a') \in A \times A: a+a' \in B \}| = \sum_{b \in B} r(b)$$
  where $r: G \to \N$ is the counting function
  $$ r(b) := |\{ (a,a') \in A \times A: a+a' = b \}|.$$
  From double counting we have
  $$ \sum_{b \in G} r(b)^2 = |G|^3 E(A).$$
  The Cauchy--Schwarz inequality then gives
  $$ \Bigl( \sum_{b \in B} r(b) \Bigr)^2 \leq |B| \sum_{b \in B} r(b)^2
     \leq |B| \sum_{b \in G} r(b)^2 = |B|\, |G|^3 E(A),$$
  and the claim follows on dividing both sides by $|G|^4$.
\end{proof}

\begin{lemma}[Balog--Szemer\'edi--Gowers lemma]
  \label{bsg}
  \lean{BSG_self}
  \uses{energy-def}
  \leanok

  Let $G$ be a finite abelian group, and let $A$ be a non-empty subset of $G$
  with $E(A) \geq \frac{1}{K} \left( \frac{|A|}{|G|} \right)^3$ for some $K \geq 1$.
  Then there is a subset $A'$ of $A$ with $|A'| \geq |A| / (C_1 K^{C_2})$ and
  $|A'-A'| \leq C_3 K^{C_4} |A|$, where
  $$C_1 = 2^4, C_2 = 1, C_3 = 2^{10}, C_4 = 5.$$
  The Lean statement gives the two conclusions as inequalities between densities; these are
  equivalent to the ones above, the common factor $|G|$ cancelling.
\end{lemma}
\begin{proof}
  \leanok

  See \url{https://terrytao.files.wordpress.com/2024/01/simplebsg.pdf}.
\end{proof}


\chapter{Homomorphism version of PFR}

In this section, $G, G'$ are finite abelian $2$-groups.

\begin{lemma}[Hahn-Banach type theorem]\label{hb-thm}\lean{hahn_banach}\leanok  Let $H_0$ be a subgroup of $G$.  Then every homomorphism $\phi: H_0 \to G'$ can be extended to a homomorphism $\tilde \phi: G \to G'$.
\end{lemma}

\begin{proof}\leanok  By induction it suffices to treat the case where $H_0$ has index $2$ in $G$, but then the extension can be constructed by hand.
\end{proof}

\begin{lemma}[Goursat type theorem]\label{goursat}\lean{goursat}\leanok  Let $H$ be a subgroup of $G \times G'$.  Then there exists a subgroup $H_0$ of $G$, a subgroup $H_1$ of $G'$, and a homomorphism $\phi: G \to G'$ such that
$$ H := \{ (x, \phi(x) + y): x \in H_0, y \in H_1 \}.$$
In particular, $|H| = |H_0| |H_1|$.
\end{lemma}

\begin{proof}\uses{hb-thm}\leanok We can take $H_0$ to be the projection of $H$ to $G$, and $H_1$ to be the slice $H_1 := \{ y: (0,y) \in H \}$.  One can construct $\phi$ on $H_0$ one generator at a time by the greedy algorithm, and then extend to $G$ by \Cref{hb-thm}.  The cardinality bound is clear from direct counting.
\end{proof}

\begin{theorem}[Approximate homomorphism form of PFR]\label{approx-hom-pfr}\lean{approx_hom_pfr}\leanok Let $G,G'$ be finite abelian $2$-groups.
  Let $f: G \to G'$ be a function, and suppose that there are at least $|G|^2 / K$ pairs $(x,y) \in G^2$ such that
$$ f(x+y) = f(x) + f(y).$$
Then there exists a homomorphism $\phi: G \to G'$ and a constant $c \in G'$
such that $f(x) = \phi(x)+c$ for at least $|G| / (2 ^ {144} * K ^ {122})$
values of $x \in G$.
\end{theorem}

\begin{proof}\uses{goursat, cs-bound, bsg, pfr-9-aux'}\leanok Consider the graph $A \subset G \times G'$ defined by
$$ A := \{ (x,f(x)): x \in G \}.$$
Clearly, $|A| = |G|$.  By hypothesis, we have $a+a' \in A$ for at least
$|A|^2/K$ pairs $(a,a') \in A^2$. By \Cref{cs-bound}, this implies that
$E(A) \geq \frac{1}{K^2} \left( \frac{|A|}{|G \times G'|} \right)^3$.  Applying \Cref{bsg}, we conclude that there exists a subset
$A' \subset A$ with $|A'| \geq |A|/C_1 K^{2C_2}$ and $|A'+A'| \leq C_1C_3
K^{2(C_2+C_4)} |A'|$. Applying \Cref{pfr-9-aux'}, we may find a subspace $H
\subset G \times G'$ such that $|H| / |A'| \in [L^{-8}, L^{8}]$ and a subset
$c$ of cardinality at most $L^5 |A'|^{1/2} / |H|^{1/2}$ such that $A'
\subseteq c + H$, where $L =  C_1C_3 K^{2(C_2+C_4)}$. If we let $H_0,H_1$ be
as in \Cref{goursat}, this implies on taking projections the projection of
$A'$ to $G$ is covered by at most $|c|$ translates of $H_0$.  This implies
that
$$ |c| |H_0| \geq |A'|;$$
since $|H_0| |H_1| = |H|$, we conclude that
$$ |H_1| \leq |c| |H|/|A'|.$$

By hypothesis, $A'$ is covered by at most $|c|$ translates of $H$, and hence
by at most $|c| |H_1|$ translates of $\{ (x,\phi(x)): x \in G \}$.  As $\phi$
is a homomorphism, each such translate can be written in the form $\{
(x,\phi(x)+c): x \in G \}$ for some $c \in G'$. The number of translates is
bounded by
$$
|c|^2 \frac{|H|}{|A'|} \leq \left(L^5 \frac{|A'|^{1/2}}{|H|^{1/2}}\right)^2 \frac{|H|}{|A'|} = L^{10}.
$$

By the pigeonhole principle, one of these translates must then contain at
least $|A'|/L^{10} \geq |G| / (C_1C_3 K^{2(C_2+C_4)})^{10} (C_1 K^{2C_2})$
elements of $A'$ (and hence of $A$), and the claim follows.
\end{proof}

\end{document}
